% !TEX root = ../thesis.tex

\chapter{Požiadavky na systém}

Táto kapitola vymedzuje základné požiadavky na navrhovaný informačný systém pre distribúciu pitnej a úžitkovej vody vrátane evidencie revízií filtračných zariadení. Zameriava sa na opis problému, motiváciu riešenia a formuláciu hlavných cieľov systému tak, aby bolo zrejmé, aké funkcie má riešenie poskytovať a aké prínosy sa od jeho nasadenia očakávajú. Požiadavky uvedené v tejto kapitole zároveň vytvárajú východisko pre návrh architektúry systému, evidenciu údajov, podporu monitorovania a následné vyhodnotenie navrhnutého riešenia.

\section{Problém a motivácia}

Vodárenské organizácie spravujú veľké množstvo aktív a prevádzkových dát; často využívajú viacero informačných systémov a dáta bývajú rozptýlené medzi oddeleniami, čo sťažuje ich spojenie do použiteľnej informácie \cite{CarricoFerreira2021}. Zároveň legislatíva Európskej únie aj Slovenskej republiky kladie dôraz na systematický monitoring a riadenie rizík v celom reťazci zásobovania pitnou vodou \cite{EUDWD2020,SKVyhlaska912023} a monitoring môže zahŕňať veľmi veľký rozsah ukazovateľov \cite{UVZPitnaVodaMapy}. Motiváciou práce je vytvoriť jednotnú a auditovateľnú evidenciu revízií filtračných zariadení, ktorá bude zahŕňať termíny, výsledky, protokoly aj nápravné opatrenia, s cieľom znížiť riziko zmeškaných činností a zrýchliť orientáciu pri kontrole alebo incidente.

\section{Hlavné ciele systému}

\begin{enumerate}[label=\arabic*.]
\item Systém musí poskytovať jednotný register prvkov distribúcie vody a filtračných zariadení vrátane identifikácie, umiestnenia, parametrov a stavu.
\item Systém musí umožniť plánovať a evidovať revízie filtračných zariadení vrátane termínu, vykonávateľa, výsledku, protokolu alebo prílohy a nápravného opatrenia s históriou zmien.
\item Systém musí podporiť evidenciu udalostí, ako sú porucha, zásah alebo zistenie, a ich prepojenie na riziká a monitoring v duchu rizikovo orientovaného prístupu \cite{EUDWD2020,SKVyhlaska912023}.
\item Systém musí automaticky zobrazovať zoznam revízií pred termínom a po termíne a generovať upozornenia pre zodpovedné osoby.
\item Systém musí generovať exportovateľné prehľady, napríklad vo formáte PDF alebo CSV, o revíziách, zisteniach a nápravných opatreniach pre manažérsku kontrolu a audit.
\end{enumerate}

